% !TeX spellcheck = es_ES
\documentclass[letterpaper,11pt]{article}
\usepackage[utf8]{inputenc}
\usepackage[T1]{fontenc}
\usepackage[spanish]{babel}
\decimalpoint % separador decimal con punto, como en las tablas de la fuente original
\usepackage{lmodern,microtype}
\usepackage{amsmath,amssymb,amsfonts}
\usepackage{geometry}
\geometry{margin=2.54cm}
\usepackage{setspace}
\setstretch{1.15}
\usepackage{enumitem}
\usepackage{booktabs}
\usepackage{color}
\usepackage{hyperref}
\definecolor{ntr}{rgb}{0.8,0,0}
\hypersetup{colorlinks=true,linkcolor=ntr,citecolor=ntr,urlcolor=black}

% Portada (Luis Cabral): titulo \LARGE sans-serif alineado a la izquierda, sin numero de pagina
\makeatletter
    \renewcommand{\@maketitle}{%
    \newpage\parindent0in\null\vskip2em%
    {\LARGE\textbf{\textsf{\@title}}\par\vskip2em}%
    \@author\par\vskip2em%
    \@date%
    \par
    \thispagestyle{empty}\setcounter{page}{1}%
    }%
\makeatother

\DeclareMathOperator{\Var}{Var}
\DeclareMathOperator{\Cov}{Cov}
\DeclareMathOperator{\Corr}{Corr}
\DeclareMathOperator{\E}{\mathbb{E}}
\DeclareMathOperator{\N}{Normal}
\DeclareMathOperator{\LP}{\mathbb{L}}
\newcommand\independent{\protect\mathpalette{\protect\independenT}{\perp}}
\def\independenT#1#2{\mathrel{\rlap{$#1#2$}\mkern2mu{#1#2}}}
% Teoria asintotica
\newcommand{\SR}{\mathbb{R}}
\DeclareMathOperator{\plim}{plim}
\DeclareMathOperator{\Avar}{AVar}
\newcommand{\pto}{\stackrel{p}{\longrightarrow}}
\newcommand{\dto}{\stackrel{d}{\longrightarrow}}
\newcommand{\adistr}{\stackrel{a}{\sim}}
\newcommand{\epsi}{\varepsilon}


\title{Análisis Econométrico \\ Ayudantía 10}
\author{Magíster en Economía, Universidad Alberto Hurtado}
\date{}

% Funciones de R sugeridas: entre corchetes, alineadas a la derecha; si no caben en la
% linea pasan enteras a la siguiente.
\newcommand{\fn}[1]{{\unskip\nobreak\hfill\penalty50\hskip1em\null\nobreak\hfill
    \mbox{[#1]}\parfillskip=0pt\finalhyphendemerits=0\par}}

\begin{document}
\maketitle

\noindent Temas: cómo estimar por variables instrumentales en R. Primera etapa, forma
reducida, mínimos cuadrados indirectos y estimador de Wald; 2SLS ``a mano'' y por qué sus
errores estándar están mal; \texttt{ivreg} con errores estándar robustos; variables
instrumentales con controles y con muchas dummies; y un modelo con varios instrumentos, con
el test de relevancia conjunta y los tests de Sargan y de Wu--Hausman.

\medskip
\noindent Toda la ayudantía es una demostración en R: el ayudante corre el código en
pantalla, así que no hace falta traer computador. Las partes 1(e)--(f) y el ejercicio 2
adelantan el modelo con controles y con varios instrumentos, que se ve en clase el lunes 16
de noviembre: acá el foco es cómo se estima en R.

\section*{Si quiere replicarlo en su computador}

\noindent El script completo, \texttt{TA10.R}, y los datos están en el repositorio del curso,
en \href{https://github.com/ramirode/econometria_mae/tree/main/TA/TA10}{\texttt{TA/TA10}}. Los
pasos son los mismos de la Ayudantía 5: cree una carpeta \texttt{TA10} con las subcarpetas
\texttt{code} (para \texttt{TA10.R}) e \texttt{input} (para \texttt{divorcio.rds} y
\texttt{bancos.rds}), y cambie la línea \texttt{dir\_proyecto <- "..."} por la ruta a su
carpeta \texttt{TA10}. Paquetes: \texttt{sandwich}, \texttt{lmtest}, \texttt{car},
\texttt{ivreg} y \texttt{fixest}.

\section*{Ejercicios}

\begin{enumerate}[itemsep=0.8em]

%% ---------------------------------------------------------------- 1
\item \emph{Divorcio y situación económica de la mujer.} Bedard y Deschênes
    (\textit{Journal of Human Resources}, 2005) estudian el efecto del divorcio sobre la
    situación económica de la mujer. Con los datos \texttt{divorcio.rds} (Censo de EE.UU. de
    1980: 576{,}017 mujeres de 21 a 40 años, alguna vez casadas y con al menos un hijo),
    considere
    \begin{align*}
    ingreso_i=\beta_0+\beta_1\, divorcio_i+u_i,
    \end{align*}
    donde $ingreso$ es el ingreso del hogar por persona (en dólares) y $divorcio$ es una
    dummy igual a 1 si el primer matrimonio de la mujer terminó. Los autores proponen como
    instrumento $nina$, una dummy igual a 1 si el primer hijo es mujer: en EE.UU., los
    matrimonios cuyo primer hijo es mujer se divorcian algo más.
    \begin{enumerate}
        \item Calcule la media y la desviación estándar de $ingreso$, $divorcio$ y $nina$.
              \fn{\texttt{readRDS}, \texttt{mean}, \texttt{sd}}
        \item Estime el modelo por MCO con errores estándar robustos. ¿Cree que
              $\hat\beta_{1}^{MCO}$ mide el efecto causal del divorcio? ¿En qué dirección
              esperaría el sesgo?
              \fn{\texttt{lm}, \texttt{vcovHC}, \texttt{coeftest}}
        \item Estime la primera etapa ($divorcio$ sobre $nina$) y la forma reducida
              ($ingreso$ sobre $nina$). Interprete los coeficientes de $nina$. Calcule el
              estimador de mínimos cuadrados indirectos como el cociente de ambos
              coeficientes, y el estimador de Wald con diferencias de medias.
              \fn{\texttt{lm}, \texttt{coef}, \texttt{tapply}}
        \item Estime por 2SLS ``a mano'': guarde el valor predicho de la primera etapa y
              regrese $ingreso$ sobre él por MCO. Luego estime con \texttt{ivreg}, con
              errores estándar robustos. Compare los coeficientes y los errores estándar.
              ¿Cuál error estándar es el correcto? ¿Es el efecto estadísticamente
              significativo? Compare con MCO.
              \fn{\texttt{fitted}, \texttt{lm}, \texttt{ivreg}, \texttt{vcovHC}}
        \item Agregue como controles la edad, la edad al primer matrimonio, la edad al
              nacer el primer hijo y los años de educación (cada una con su cuadrado), y la
              dummy $metro$. Estime por MCO y por variables instrumentales con
              \texttt{feols} del paquete \texttt{fixest} (con 576 mil observaciones y
              tantos regresores, \texttt{lm} e \texttt{ivreg} usan mucha memoria).
              Contraste la relevancia del instrumento en la primera etapa con controles.
              ¿Cambian mucho los resultados?
              \fn{\texttt{feols}, \texttt{coeftable}, \texttt{fitstat}}
        \item Agregue dummies de estado de residencia y de lugar de nacimiento (unas 100
              dummies) y estime de nuevo con \texttt{feols}, que trata las dummies como
              efectos fijos. Compare con (e).
              \fn{\texttt{feols}, \texttt{coeftable}, \texttt{fitstat}}
    \end{enumerate}

%% ---------------------------------------------------------------- 2
\item \emph{Bancos y desarrollo.} Pascali (\textit{Review of Economics and Statistics},
    2016) estudia si el desarrollo financiero de las ciudades italianas afecta su ingreso
    per cápita. Con los datos \texttt{bancos.rds} (municipios del centro y norte de Italia,
    2002 y 2003), considere
    \begin{align*}
    pib_{it}=\beta_0+\beta_1\, credito_{it}+x_{it}'\gamma+u_{it},
    \end{align*}
    donde $pib$ es el valor agregado per cápita (miles de euros), $credito$ es el crédito
    privado como porcentaje del PIB y $x$ incluye características geográficas del municipio
    (altitud, clase sísmica, acceso al mar y a un río) y dummies de provincia y de año. El
    instrumento propuesto es $jew$, una dummy igual a 1 si en el municipio hubo una
    comunidad judía en el siglo XV: en el Renacimiento, para competir con los prestamistas
    judíos se fundaron instituciones de crédito locales (los \textit{Monti di Pietà}), muchas
    de las cuales se convirtieron después en bancos.
    \begin{enumerate}
        \item Estime el modelo por MCO con errores estándar robustos. ¿Cree que mide un
              efecto causal?
              \fn{\texttt{lm}, \texttt{factor}, \texttt{coeftest}}
        \item Estime la primera etapa con $jew$ como instrumento y contraste su relevancia.
              ¿Qué condiciones debe cumplir $jew$ para ser un instrumento válido? ¿Cuál se
              puede contrastar?
              \fn{\texttt{lm}, \texttt{linearHypothesis}}
        \item Estime el modelo por variables instrumentales con $jew$. Compare con MCO.
              \fn{\texttt{ivreg}, \texttt{vcovHC}}
        \item Use ahora tres instrumentos: dummies de comunidad judía pequeña ($jew\_peq$),
              mediana ($jew\_med$) y grande ($jew\_gde$). Estime la primera etapa y
              contraste la relevancia conjunta de los tres instrumentos. Estime por 2SLS y
              reporte los tests que entrega \texttt{summary(\dots, diagnostics = TRUE)}:
              instrumentos débiles, Wu--Hausman y Sargan. ¿Qué hipótesis contrasta cada uno
              y qué concluye?
              \fn{\texttt{linearHypothesis}, \texttt{ivreg}, \texttt{summary}}
    \end{enumerate}

\end{enumerate}

\end{document}
